\chapter{Conclusion}

This thesis investigated whether a server can mitigate free riding using only returned model updates, without trusting client-reported accuracy, loss, training time, or resource use and without training an auxiliary detector. SWT-CP v8 addresses this objective through secret trap models, two-sided norm and layer-profile evidence, a candidate/suspicion state machine, repeated probing, rotating anchors, cumulative penalties, and rehabilitation.

The IID evaluation is encouraging: every tested free rider was removed across attacker shares of 30\% and 40\% and seeds 42 and 43. Seed 43 produced only one honest removal at 30\% and none at 40\%. FR1--FR3 also retained complete recall under non-IID partitioning. These results support the use of unpredictable server interventions and longitudinal evidence for lightweight detection of replay, random-update, and historical-average behavior.

The evaluation also revealed two limits that are as important as the successes. First, non-IID false positives concentrate on particular honest clients and recur across attacks, showing that joint statistical evidence is still sensitive to legitimate client heterogeneity. Second, FR4 exhibits a reproducible contamination transition: both non-IID 30\% runs removed all attackers, whereas both 40\% runs removed none. Undetected FR4 clients entered the anchor roster, and malicious-majority panels directly contributed to an honest removal. Rotation is therefore not equivalent to trust establishment.

The immediate research priority is to replace negative anchor eligibility (``never flagged'') with positive, repeated validation, and to use differential responses from the same client across independent traps so that each client partly serves as its own non-IID control. A third seed and intermediate 35\% FR4 experiment should confirm the transition. Broader validation should include CIFAR-10, additional Dirichlet concentrations, client dropout, on--off and partial-training free riders, collusion, defense-aware adaptation, and comparisons with implemented baselines. A secure protocol is also needed if per-client auditing is to coexist with secure aggregation.

SWT-CP should thus be understood as a computationally simple and empirically strong detector within a defined operating region, not as proof that a client trained. Its main contribution is both constructive and diagnostic: it demonstrates where secret statistical traps work, and identifies the precise combination of heterogeneity, adaptive fabrication, and contaminated references under which they fail.
